\begin{textsample}{POS Dim 4 – ap – Score 4.00 – ap/ap\_em\_14.txt}  \label{ex:f4_pos_240}
A EDUCAÇÃO ESCOLAR \textbf{INDÍGENA} NO ESTADO DO AMAPÁ No Estado do Amapá, a educação escolar \textbf{indígena} é gerenciada pela Secretaria de Estado da Educação ( SEED ), por meio do Núcleo de Educação \textbf{Indígena} ( NEI ), criado como uma seção na Divisão de Ensino de Primeiro Grau, através da Portaria No 966 de 27 de dezembro de 1990, com a incumbência principal de planejar e implementar a política de educação escolar \textbf{indígena} do Estado, em consonância com as deliberações definidas em Assembleias \textbf{Indígenas}.
Dentre as finalidades do NEI estão : contribuir para a definição dos parâmetros da política de educação escolar \textbf{indígena}, garantindo a valorização das culturas, línguas e tradições dos povos \textbf{indígenas}, respeitando as peculiaridades e demandas de cada comunidade ; propor, articular, apoiar, assessorar, acompanhar e avaliar a execução da política de educação escolar \textbf{indígena} intercultural, bilíngue, específica e diferenciada, conforme preceituam a Constituição Federal, a LDB Lei n.o 9.304/96 e as normatizações dos Conselhos Nacional e Estadual de Educação.
O NEI tem como atribuições : formular, coordenar e acompanhar as ações voltadas à política da educação escolar \textbf{indígena} ; encaminhar ao Conselho Estadual de Educação pedidos de autorização para funcionamento de Escola \textbf{Indígena} ; apresentar sugestões para melhoria da qualidade das escolas \textbf{indígenas} ; diagnosticar junto às comunidades \textbf{indígenas} as necessidades de recursos humanos, físicos e didático-pedagógicos nas escolas das aldeias ; promover a formação continuada de professores \textbf{indígenas}, levando -se em conta a língua e a cultura de cada etnia, como também os saberes \textbf{indígenas} ; estimular a contratação de professores e funcionários \textbf{indígenas}, indicados pelas comunidades ; solicitar relatórios aos professores \textbf{indígenas} sobre o desempenho da escola e dos estudantes ; e acompanhar, avaliar e emitir parecer sobre o funcionamento das Escolas Estaduais \textbf{Indígenas}.
Tais finalidades e atribuições, para serem efetivamente desenvolvidas, requer a implementação de políticas públicas a nível Federal e Estadual voltadas para o que já está garantido legalmente aos povos \textbf{indígenas}.
Atualmente, o NEI funciona em duas salas, uma ocupada pela gerência e outra onde os professores e diretores \textbf{indígenas} e não \textbf{indígenas} são recebidos para : orientações referentes à gestão administrativa, pedagógica e financeira das escolas \textbf{indígenas} e ; o assessoramento pedagógico.
No mesmo espaço físico também são recebidas à comunidade e as lideranças das diversas etnias.
Administrativamente, o NEI é composto por uma gerencia e três unidades : pedagógica, antropológica e linguística que são cargos remunerados indicados e homologados pelo Governo do Estado.
Cada Unidade tem suas competências específicas.
Compete a unidade pedagógica : diagnosticar a educação nas comunidades \textbf{indígenas}, identificando os processos próprios de aprendizagem e o papel da escolarização para implementar ações educacionais, de acordo com as demandas existentes ; discutir, planejar e coordenar a elaboração de propostas curriculares e pedagógicas para a educação escolar \textbf{indígena}, garantindo as especificidades culturais, linguísticas e pedagógicas ; supervisionar e orientar a ação pedagógica dos professores que atuam na educação escolar nas áreas \textbf{indígenas} ; planejar e orientar a confecção de matérias didáticos e pedagógicos, a partir de propostas curriculares específicas.
A unidade antropológica tem como competências : avaliar o impacto da escolarização nas áreas \textbf{indígenas} em relação à identidade étnica das culturas \textbf{indígenas} ; diagnosticar as relações de interculturalidade entre as sociedades \textbf{indígenas}, para que as ações educacionais atendam os princípios de alteridade e manutenção cultural ; elaborar planos, programas e projetos de educação escolar, voltados a ações antropológicas, garantindo as especificidades culturais de cada sociedade \textbf{indígena}.
Enquanto que a unidade linguística compete : identificar a situação sociolinguística nas áreas \textbf{indígenas} do Estado, em relação ao padrão de uso interativo das línguas \textbf{indígenas}, monolinguismo e grau de bilinguismo para implementar ações educacionais que expressem as demandas das sociedades \textbf{indígenas} ; propor uma política linguística para as ações educacionais, que visem a valorização e manutenção das línguas \textbf{indígenas} ; elaborar e acompanhar ações de planejamento linguístico para o desenvolvimento de competência intercultural, através da alfabetização em língua materna e aprendizado de língua de interação nacional.
O núcleo coordena a educação escolar de nove etnias distribuídas em terras \textbf{indígenas} situadas no Estado do Amapá e \textbf{norte} do Pará, de acordo com a tabela abaixo : Estas etnias foram organizadas pelo Núcleo de Educação \textbf{Indígena} em três regiões geográficas para o assessoramento administrativo, financeiro, antropológico, linguístico e pedagógico : Oiapoque, Parque do Tumucumaque e Pedra Branca do Amapari.
A distribuição foi feita de acordo com as suas respectivas terras \textbf{indígenas}.
As escolas estaduais que atendem os estudantes estão localizadas nas terras \textbf{indígenas}, no que concerne ao ensino médio, é ofertado em todas as terras \textbf{indígenas} do Estado do Amapá pelo Sistema de Organização Modular \textbf{Indígena} ( SOMEI ) de forma regular ou por meio da Educação de Jovens e Adultos ( EJA ).
Em Pedra Branca do Amapari, o ensino médio é ofertado na Aldeia Aramirã, considerada a central do povo Wajãpi.
Nas terras \textbf{indígenas} de Oiapoque nas Aldeias : Manga, Santa Izabel, Espírito Santo, Açaizal, Estrela do povo Karipuna ; na aldeia Kumarumã do povo Galibi Marworno ; nas aldeias Kumenê e Flexa do povo Palikur ; na aldeia São José do povo Galibi Kalinã.
Em relação ao Parque do Tumucumaque o ensino médio é ofertado na aldeia Bona considerada aldeia central dos povos Apalai, Wayana e Tiryo.

% matched lemmas: indígena, norte
\end{textsample}
